% chap03.tex 第三章 控制系统的时域分析（大气排版版）
\chapter{控制系统的时域分析}

% ===== 局部排版优化（仅影响本章，使版面更舒展）=====
\setlength{\parskip}{0.8ex plus 0.3ex minus 0.2ex}   % 段落间距增加
\setlist{
    leftmargin=2.5em,          % 列表左缩进增大
    itemsep=0.6em,             % 条目垂直间距
    parsep=0.3em,              % 条目内段落间距
    topsep=0.8em,              % 列表与上下文间距
    partopsep=0.3em
}
\renewcommand{\arraystretch}{1.5}      % 表格行高增加
\setlength{\tabcolsep}{10pt}           % 表格列间距增加
\addtolength{\abovedisplayskip}{2pt}   % 公式前间距微增
\addtolength{\belowdisplayskip}{2pt}   % 公式后间距微增
% ========================================

\bigskip  % 章节标题后额外留白

\section{本章考核要点}
\begin{enumerate}
    \item 典型输入信号与时域响应分类（动态/稳态）；
    \item 一阶、二阶系统的单位阶跃响应与时域性能指标计算；
    \item 高阶系统的主导极点简化方法；
    \item 劳斯判据稳定性分析；
    \item 稳态误差计算（静态误差系数法、终值定理法）。
\end{enumerate}
\noindent\textbf{补充：}时域法是最基本的分析方法，是复域法、频域法的基础，直接在时间域对系统进行分析和设计，直观准确，但求解解析解过程较繁琐。

\medskip

\section{系统时间响应与性能指标概述}
\subsection{课程任务与体系结构}
控制系统的分析与综合构成闭环逻辑：
\begin{enumerate}
    \item \textbf{系统分析}：由系统模型计算控制系统的各项性能指标，判断其是否满足设计要求，包括时域法、复域法、频域法三类方法；
    \item \textbf{系统综合}：根据性能指标要求，寻求改进系统性能、使其满足设计规范的校正方法。
\end{enumerate}

\subsection{本章主要内容}
\begin{enumerate}
    \item 3.1 系统时间响应的性能指标；
    \item 3.2 一阶系统的时域分析；
    \item 3.3 二阶系统的时域分析；
    \item 3.4 高阶系统的时域分析；
    \item 3.5 线性系统的稳定性分析；
    \item 3.6 线性系统的稳态误差计算。
\end{enumerate}

\subsection{系统时间响应的求解步骤}
分析系统性能的通用步骤：
\begin{enumerate}
    \item 给出控制系统的闭环传递函数 $\Phi(s)$；
    \item 用 $Y(s) = \Phi(s)R(s)$ 求出系统输出的拉氏变换 $Y(s)$，再通过拉普拉斯逆变换求 $y(t)$；
    \item 根据 $y(t)$ 分析系统动态与稳态性能；
    \item 建立系统性能与传递函数参数之间的对应关系。
\end{enumerate}
\[
\text{数学模型} \xrightarrow{\text{典型输入下求解微分方程}} \text{时域响应} \rightarrow \text{性能指标}
\]

\subsection{典型输入信号}
控制系统常用的典型测试信号，用于统一标准下的性能对比：
\begin{table}[htbp]
\centering
\caption{典型输入信号及其拉普拉斯变换}
\begin{tabular}{ccccc}
\toprule
信号名称 & 时域表达式 $f(t)$ & 拉氏变换 $F(s)$ & 典型应用场景 \\
\midrule
单位脉冲 & $\delta(t)$ & $1$ & 冲击信号、电脉冲测试 \\
单位阶跃 & $1(t) = \begin{cases} 1, & t\ge0 \\ 0, & t<0 \end{cases}$ & $\dfrac{1}{s}$ & 开关量、阶跃扰动测试 \\
单位斜坡 & $t\cdot1(t) = \begin{cases} t, & t\ge0 \\ 0, & t<0 \end{cases}$ & $\dfrac{1}{s^2}$ & 等速跟踪、恒速输入测试 \\
单位加速度 & $\dfrac{t^2}{2}\cdot1(t) = \begin{cases} t^2/2, & t\ge0 \\ 0, & t<0 \end{cases}$ & $\dfrac{1}{s^3}$ & 恒加速跟踪测试 \\
\bottomrule
\end{tabular}
\end{table}
信号之间的时域关系：单位阶跃是单位脉冲的积分，单位斜坡是单位阶跃的积分，单位加速度是单位斜坡的积分，对应复域关系为乘以 $1/s$。

\subsection{动态过程与稳态过程}
系统响应按时间阶段分为两类：
\begin{enumerate}
    \item \textbf{动态过程（过渡过程/瞬态过程）}：典型输入作用下，系统输出从初始状态到最终状态的响应过程。稳定系统的动态过程是收敛的，不稳定系统的动态过程是发散或等幅振荡的。
    \item \textbf{稳态过程（稳态响应）}：典型输入作用下，当 $t\to\infty$ 时的系统输出，表征系统输出最终复现输入量的程度，用稳态性能指标描述。
\end{enumerate}

\subsection{动态性能与稳态性能}
控制系统的基本性能要求可概括为“稳、准、快”：
\begin{itemize}
    \item \textbf{稳（稳定性）}：系统受扰动影响后能回到原来的平衡位置，是系统正常工作的基本前提；
    \item \textbf{准（稳态性能）}：稳态输出与理想输出之间的误差（稳态误差）要小，反映系统的控制精度；
    \item \textbf{快（动态性能）}：阶跃响应的过渡过程要平稳、迅速，反映系统的动态响应速度与平稳性。
\end{itemize}
稳态误差是描述系统稳态性能的核心指标，通常在阶跃、斜坡或加速度信号作用下计算。若 $t\to\infty$ 时系统输出不等于输入或输入的确定函数，则存在稳态误差，其大小反映了系统的控制精度和抗扰动能力。

\subsection{线性系统时域性能指标（单位阶跃响应）}
仅对稳定系统定义以下时域性能指标：
\begin{itemize}
    \item \textbf{延迟时间 $t_d$}：阶跃响应第一次达到终值的 $50\%$ 所需的时间；
    \item \textbf{上升时间 $t_r$}：阶跃响应从终值的 $10\%$ 上升到终值的 $90\%$ 所需的时间；有振荡时，也可定义为从 $0$ 到第一次达到终值所需的时间；
    \item \textbf{峰值时间 $t_p$}：阶跃响应越过终值达到第一个峰值所需的时间；
    \item \textbf{超调量 $\sigma\%$}：峰值超出终值的百分比：
    \[
    \sigma\% = \frac{y(t_p) - y(\infty)}{y(\infty)} \times 100\%
    \]
    \item \textbf{调节时间 $t_s$}：阶跃响应到达并保持在终值的 $5\%$ 或 $2\%$ 误差带内所需的最短时间，满足 $|y(t)-y(\infty)| \le \Delta y(\infty),\ t\le t_s$；
    \item \textbf{振荡次数 $N$}：在 $t\le t_s$ 时间内，单位阶跃响应穿越稳态值 $y(\infty)$ 次数的一半。
\end{itemize}

\bigskip

\section{一阶系统的时域分析}
一阶系统标准闭环传递函数：
\[
\Phi(s)=\frac{C(s)}{R(s)}=\frac{1}{Ts+1}
\]
$T$ 为一阶系统的时间常数，直接决定系统动态响应速度与稳态跟踪能力。

\subsection{一阶系统单位阶跃响应回顾}
输入为单位阶跃信号 $R(s)=\frac{1}{s}$，输出：
\[
Y(s)=\frac{1}{Ts+1}\cdot\frac{1}{s}
\]
拉氏反变换得时域响应：
\[
y(t)=1-e^{-\frac{t}{T}},\quad t\ge 0
\]
关键特征：$y(T)=0.632$、$y(2T)=0.865$、$y(3T)=0.95$、$y(4T)=0.98$；\textbf{时间常数 $T$ 越小，响应上升速度越快，调节时间越短}。

思考：4条阶跃响应曲线时间常数分别为 $0.5,1,2,5$，曲线上升速度从快到慢对应 $T=0.5>1>2>5$，最陡曲线对应 $T=0.5$，最平缓曲线对应 $T=5$。

\subsection{一阶系统单位脉冲响应}
输入为单位脉冲信号 $r(t)=\delta(t),\ R(s)=1$，输出：
\[
Y(s)=\frac{1}{Ts+1}
\]
拉氏反变换得：
\[
y(t)=\frac{1}{T}e^{-\frac{t}{T}},\quad t\ge 0
\]
关键特征参数：
\[
\begin{cases}
y(0)=\displaystyle\frac{1}{T}\\
y(\infty)=0\\
y'(0)=-\displaystyle\frac{1}{T^2}\\
y(T)=\displaystyle\frac{0.368}{T}
\end{cases}
\]
响应特点：
\begin{enumerate}
    \item $y(t)$ 是单调连续下降的指数曲线；
    \item 终值为 $0$，脉冲输入下稳态误差为 $0$；
    \item 动态性能完全由时间常数 $T$ 决定；
    \item 脉冲响应做拉氏变换可直接得到系统传递函数，完整表征系统特性。
\end{enumerate}

\subsection{一阶系统单位斜坡响应}
输入为单位斜坡信号 $r(t)=t,\ R(s)=\frac{1}{s^2}$，输出：
\[
Y(s)=\frac{1}{Ts+1}\cdot\frac{1}{s^2}
\]
拉氏反变换得：
\[
y(t)=t-T\left(1-e^{-\frac{t}{T}}\right),\quad t\ge 0
\]
响应特点：
\begin{enumerate}
    \item $t\to\infty$ 时，$y(t)\to t-T$；
    \item 跟踪斜坡信号存在稳态误差，稳态误差数值等于时间常数 $T$；
    \item 减小 $T$ 既能加快响应速度，也能减小斜坡跟踪的稳态误差。
\end{enumerate}

\subsection{一阶系统单位匀加速度响应}
输入为单位匀加速度信号 $r(t)=\frac12 t^2,\ R(s)=\frac{1}{s^3}$，输出：
\[
Y(s)=\frac{1}{Ts+1}\cdot\frac{1}{s^3}
\]
拉氏反变换得：
\[
y(t)=\frac12 t^2-Tt+T^2\left(1-e^{-\frac{t}{T}}\right),\quad t\ge 0
\]
响应特点：
\begin{enumerate}
    \item 跟踪误差会随时间推移持续增大；
    \item 一阶系统无法实现对匀加速度信号的有效跟踪。
\end{enumerate}

\subsection{一阶系统时域分析小结}
\begin{table}[htbp]
\centering
\caption{一阶系统对典型输入信号的响应汇总}
\begin{tabular}{ccccc}
\toprule
输入 $r(t)$ & $R(s)$ & $C(s)=\Phi(s)R(s)$ & 时域响应 $c(t)$ \\
\midrule
$\delta(t)$ & $1$ & $\displaystyle\frac{1}{Ts+1}$ & $\displaystyle k(t)=\frac{1}{T}e^{-\frac{t}{T}}$ \\
$1(t)$ & $\displaystyle\frac{1}{s}$ & $\displaystyle\frac{1}{s}-\frac{1}{s+\frac{1}{T}}$ & $\displaystyle h(t)=1-e^{-\frac{t}{T}}$ \\
$t$ & $\displaystyle\frac{1}{s^2}$ & $\displaystyle\frac{1}{s^2}-\frac{T}{s}+\frac{T}{s+\frac{1}{T}}$ & $\displaystyle c(t)=t-T\left(1-e^{-\frac{t}{T}}\right)$ \\
$\frac12 t^2$ & $\displaystyle\frac{1}{s^3}$ & / & $\displaystyle c(t)=\frac12 t^2-Tt+T^2\left(1-e^{-\frac{t}{T}}\right)$ \\
\bottomrule
\end{tabular}
\end{table}
\begin{enumerate}
    \item \textbf{等价关系}：系统对输入导数的响应 = 对原输入响应的导数；系统对输入积分的响应 = 对原输入响应的积分；
    \item \textbf{动态特性}：时间常数 $T$ 越大，响应速度越慢、调节时间越长，反之响应越快；
    \item \textbf{跟踪能力}：阶跃输入无稳态误差；斜坡输入稳态误差等于 $T$；加速度输入稳态误差趋于无穷，无法跟踪。
\end{enumerate}

\bigskip

\section{二阶系统的时域分析}
\subsection{二阶系统的数学模型}
\subsubsection{标准形式定义}
标准二阶系统微分方程：
\[
\ddot{y}+2\zeta\omega_n \dot{y}+\omega_n^2 y=\omega_n^2 r
\]
闭环传递函数（首1标准型）：
\[
\Phi(s)=\frac{C(s)}{R(s)}=\frac{\omega_n^2}{s^2+2\zeta\omega_n s+\omega_n^2}
\]
尾1标准型闭环传递函数：
\[
\Phi(s)=\frac{1}{T^2 s^2+2T\zeta s+1}
\]
其中：
\begin{itemize}
    \item $\omega_n$：自然频率（或无阻尼振荡频率）；
    \item $\zeta$：阻尼比（或阻尼系数）。
\end{itemize}

单位反馈结构对应的开环传递函数：
\[
G(s)=\frac{\omega_n^2}{s(s+2\zeta\omega_n)}
\]

\subsubsection{电动伺服系统实例}
以电动伺服机构为例，通过电位计测偏差、放大器放大、直流电机驱动负载，可推导出上述标准二阶形式，完整框图可通过环节串联与反馈化简得到标准闭环表达式。

\subsubsection{非标准二阶系统标准化方法}
若已知单位反馈开环传函 $G(s)=\dfrac{K}{s(Ts+1)}$，变形拆分：
\[
G(s)=\frac{K/T}{s\left(s+\frac1T\right)}
\]
联立标准形式参数对应方程组：
\[
\begin{cases}
\omega_n^2=\displaystyle\frac{K}{T}\\[4pt]
2\zeta\omega_n=\displaystyle\frac1T
\end{cases}
\]
求解得：
\[
\begin{cases}
\omega_n=\sqrt{\displaystyle\frac{K}{T}}\\[4pt]
\zeta=\sqrt{\displaystyle\frac{1}{4TK}}
\end{cases}
\]
反向换算关系：
\[
\frac{K}{T}=\omega_n^2,\quad TK=\frac{1}{4\zeta^2}
\]

\subsubsection{标准二阶系统特征方程与特征根}
特征方程：
\[
s^2+2\zeta\omega_n s+\omega_n^2=0
\]
求解特征根：
\[
s_{1,2}=-\zeta\omega_n\pm\omega_n\sqrt{\zeta^2-1}
\]
阻尼比 $\zeta$ 的取值直接决定特征根在 $s$ 平面上的位置，分区间讨论：
\begin{enumerate}
    \item $\boldsymbol{\zeta>0}$，根实部为负（系统稳定）：
    \begin{itemize}
        \item $0<\zeta<1$：一对共轭复根（欠阻尼）；
        \item $\zeta=1$：一对相等实根（临界阻尼）；
        \item $\zeta>1$：一对不等实根（过阻尼）。
    \end{itemize}
    \item $\boldsymbol{\zeta=0}$，实部为0：纯虚共轭根 $s_{1,2}=\pm j\omega_n$（无阻尼，等幅振荡）；
    \item $\boldsymbol{\zeta<0}$，根实部为正（系统不稳定）：
    \begin{itemize}
        \item $-1<\zeta<0$：一对共轭复根（负阻尼，发散振荡）；
        \item $\zeta=-1$：一对相等正实根；
        \item $\zeta<-1$：一对不等正实根。
    \end{itemize}
\end{enumerate}

阻尼工况分类汇总：
\begin{itemize}
    \item $0<\zeta<1$：欠阻尼；
    \item $\zeta=1$：临界阻尼；
    \item $\zeta>1$：过阻尼；
    \item $\zeta=0$：无阻尼；
    \item $-1<\zeta<0$：负阻尼。
\end{itemize}

\subsection{二阶系统的单位阶跃响应}
\subsubsection{① 欠阻尼 $\boldsymbol{0<\zeta<1}$ 的情况}
共轭复数极点：
\[
s_{1,2}=-\zeta\omega_n\pm j\omega_n\sqrt{1-\zeta^2}
\]
定义有阻尼振荡频率：
\[
\omega_d=\omega_n\sqrt{1-\zeta^2}
\]

时域完整表达式：
\[
y(t)=1-e^{-\zeta\omega_n t}\left(\cos\omega_d t+\frac{\zeta}{\sqrt{1-\zeta^2}}\sin\omega_d t\right),\quad t\ge0
\]
等价合并正弦相位形式：
\[
y(t)=1-\frac{e^{-\zeta\omega_n t}}{\sqrt{1-\zeta^2}}\sin\left(\omega_d t+\theta\right),\quad t\ge0
\]
其中相位角：
\[
\theta=\arctan\frac{\sqrt{1-\zeta^2}}{\zeta}=\arccos\zeta
\]

将响应拆分为稳态分量、暂态分量：
\[
y(t)=\underbrace{1}_{\text{稳态分量}}-\underbrace{\frac{e^{-\zeta\omega_n t}}{\sqrt{1-\zeta^2}}\sin\left(\omega_d t+\theta\right)}_{\text{暂态分量}}
\]
响应包络线约束：$\displaystyle 1\pm \frac{e^{-\zeta\omega_n t}}{\sqrt{1-\zeta^2}}$。

\paragraph{欠阻尼单位阶跃响应特点}
\begin{enumerate}
    \item 暂态分量是衰减的正弦振荡；
    \item 暂态分量的衰减速度取决于衰减系数 $\zeta\omega_n$；
    \item 系统输出稳态值恒为 $1$；
    \item 阶跃信号作用下不存在稳态误差；
    \item 响应曲线被包络线约束，呈衰减振荡收敛至稳态。
\end{enumerate}

\subsubsection{② 无阻尼 $\boldsymbol{\zeta=0}$ 的情况}
闭环极点：
\[
s_{1,2}=\pm j\omega_n
\]
极点落在虚轴上，代入通用阶跃响应公式，$\zeta=0,\theta=90^\circ$：
\[
y(t)=1-\cos\left(\omega_n t\right),\quad t\ge0
\]

\paragraph{无阻尼响应特点}
\begin{enumerate}
    \item 单位阶跃响应为等幅振荡过程；
    \item 超调量 $\sigma_p=100\%$，调节时间 $t_s\to\infty$；
    \item 系统临界稳定，无法正常工作。
\end{enumerate}

\subsubsection{③ 临界阻尼 $\boldsymbol{\zeta=1}$ 的情况}
一对相同负实数极点：
\[
s_{1,2}=-\omega_n
\]
单位阶跃响应：
\[
y(t)=1-e^{-\omega_n t}\left(1+\omega_n t\right),\quad t\ge0
\]

求导：$\displaystyle \frac{dy(t)}{dt}=\omega_n^2 t e^{-\omega_n t}$，$t=0$ 时曲线斜率为 $0$。

\paragraph{临界阻尼响应特点}
\begin{enumerate}
    \item 稳态值为 $1$，无超调、单调上升；
    \item 曲线形态类似一阶系统阶跃响应，但初始斜率为0，二者存在本质区别。
\end{enumerate}

\subsubsection{④ 过阻尼 $\boldsymbol{\zeta>1}$ 的情况}
两个不相等的负实数极点：
\[
s_{1,2}=-\zeta\omega_n\pm\omega_n\sqrt{\zeta^2-1}
\]
单位阶跃响应：
\[
\begin{aligned}
y(t) &= 1-\frac{1}{2\sqrt{\zeta^2-1}\left(\zeta-\sqrt{\zeta^2-1}\right)}e^{-\left(\zeta-\sqrt{\zeta^2-1}\right)\omega_n t} \\
&\quad +\frac{1}{2\sqrt{\zeta^2-1}\left(\zeta+\sqrt{\zeta^2-1}\right)}e^{-\left(\zeta+\sqrt{\zeta^2-1}\right)\omega_n t},\quad t\ge0
\end{aligned}
\]
近似表达式（远离重极点时）：
\[
y(t) \approx 1-\frac{1}{2\sqrt{\zeta^2-1}\left(\zeta-\sqrt{\zeta^2-1}\right)}e^{-\left(\zeta-\sqrt{\zeta^2-1}\right)\omega_n t}
\]

\paragraph{过阻尼响应特点}
稳态值为 $1$，全程单调上升，无振荡、无超调。

\subsubsection{⑤ 负阻尼 $\boldsymbol{-1<\zeta<0}$ 的情况}
阶跃响应通用形式：
\[
y(t)=1-\frac{e^{-\zeta\omega_n t}}{\sqrt{1-\zeta^2}}\sin\left(\omega_d t+\theta\right),\quad t\ge0
\]
其中：
\[
\theta=\arctan\frac{\sqrt{1-\zeta^2}}{\zeta},\quad \omega_d=\omega_n\sqrt{1-\zeta^2}
\]

\paragraph{负阻尼响应特点}
暂态分量指数项 $e^{-\zeta\omega_n t}$ 随时间指数发散，输出为发散正弦振荡，系统不稳定。

\subsection{欠阻尼二阶系统的动态过程指标}
\subsubsection{① 上升时间 $t_r$}
上升时间定义：单位阶跃响应 $y(t)$ 第一次达到其稳态值 $y(\infty)$ 所需的时间；对于过阻尼系统，通常采用 $y(t)$ 从稳态值的 $10\%$ 上升到 $90\%$ 所需的时间，$t_r$ 可以反映系统的快速性。

令 $y(t_r)=1$：
\[
y(t_r) = 1 - e^{-\zeta\omega_n t_r}\sin(\omega_d t_r+\theta) = 1
\]
化简得到：
\[
\sin(\omega_d t_r+\theta)=0
\]
取第一次抵达稳态的最小正根：
\[
t_r = \frac{\pi-\theta}{\omega_d} = \frac{\pi-\theta}{\omega_n\sqrt{1-\zeta^2}}
\]
其中相位角 $\displaystyle\theta = \arctan\frac{\sqrt{1-\zeta^2}}{\zeta}$，$\theta$ 以弧度为单位。

\textbf{规律结论}：当阻尼比 $\zeta$ 一定（$\theta$ 固定）时，提高无阻尼振荡频率 $\omega_n$，可以缩短上升时间 $t_r$，系统快速性与 $\omega_n$ 成正比。

\subsubsection{② 峰值时间 $t_p$}
峰值时间定义：阶跃响应越过终值达到第一个峰值所需的时间。

对阶跃响应求导并令导数在 $t=t_p$ 处为0：
\[
\left.\frac{dy(t)}{dt}\right|_{t=t_p} = -\frac{1}{\sqrt{1-\zeta^2}}e^{-\zeta\omega_n t}\left[-\zeta\omega_n\sin(\omega_d t+\theta)+\omega_d\cos(\omega_d t+\theta)\right]\bigg|_{t=t_p}=0
\]
化简可得：
\[
\sin(\omega_d t_p+\theta-\theta)=0
\]
取第一个峰值对应的最小正根：
\[
t_p=\frac{\pi}{\omega_d}=\frac{\pi}{\omega_n\sqrt{1-\zeta^2}}
\]

\textbf{规律结论}：$\zeta$ 一定时，$\omega_n$ 增大（闭环极点离虚轴距离越远），$t_p$ 减小，响应速度变快。

\subsubsection{③ 超调量 $\sigma_p$}
将 $t_p$ 代入阶跃响应表达式：
\[
\begin{aligned}
y(t_p) &= 1-e^{-\zeta\omega_n t_p}\frac{1}{\sqrt{1-\zeta^2}}\sin(\omega_d t_p+\theta) \\
&=1-e^{-\frac{\zeta\pi}{\sqrt{1-\zeta^2}}}\frac{1}{\sqrt{1-\zeta^2}}\sin(\pi+\theta) \\
&=1+e^{-\frac{\zeta\pi}{\sqrt{1-\zeta^2}}}
\end{aligned}
\]
结合超调量定义：
\[
\sigma_p = \frac{y(t_p)-y(\infty)}{y(\infty)}\times 100\%
\]
推导出最终公式：
\[
\sigma_p = e^{-\frac{\zeta\pi}{\sqrt{1-\zeta^2}}}\times 100\%
\]

常用工程取值参考：
\[
\begin{cases}
\zeta=0,\quad \sigma_p=100\% \\
\zeta=0.4,\quad \sigma_p=25.4\% \\
\zeta=0.5,\quad \sigma_p=16.3\% \\
\zeta=\frac{\sqrt{2}}{2},\quad \sigma_p=4.3\% \\
\zeta=0.8,\quad \sigma_p=1.5\% \\
\zeta=1.0,\quad \sigma_p=0
\end{cases}
\]

\textbf{规律说明}：超调量 $\sigma_p$ 仅由阻尼比 $\zeta$ 决定，和自然频率 $\omega_n$ 无关；$\zeta$ 增大时，$\sigma_p$ 单调减小，系统平稳性提升。
由超调量反解阻尼比：
\[
\zeta=\frac{\sqrt{\ln^2 \sigma_p}}{\sqrt{\pi^2+\ln^2 \sigma_p}}=\frac{\ln \dfrac{1}{\sigma_p}}{\sqrt{\pi^2+\left(\ln \dfrac{1}{\sigma_p}\right)^2}}
\]

\subsubsection{④ 调节时间 $t_s$}
调节时间定义：单位阶跃响应达到并保持在稳态值 $\pm5\%$（或 $\pm2\%$）误差范围内所需的最短时间，满足 $|y(t)-y(\infty)|\le \Delta y(\infty),\ t\ge t_s$。

响应完整表达式：
\[
y(t) = 1 - \frac{1}{\sqrt{1-\zeta^2}}e^{-\zeta\omega_n t}\sin\left(\omega_d t + \arctan\frac{\sqrt{1-\zeta^2}}{\zeta}\right)
\]
响应的包络线为 $\displaystyle 1\pm\frac{1}{\sqrt{1-\zeta^2}}e^{-\zeta\omega_n t}$，衰减速度由 $\zeta\omega_n$ 决定，用包络线近似震荡衰减过程求解：
\[
\left|\frac{1}{\sqrt{1-\zeta^2}}e^{-\zeta\omega_n t_s}\right|\le \Delta
\]
精确解：
\[
t_s = \frac{1}{\zeta\omega_n}\ln\left(\frac{1}{\Delta\sqrt{1-\zeta^2}}\right)
\]
工程近似公式：
\begin{itemize}
    \item $2\%$ 误差带：$\displaystyle t_s=\frac{4}{\zeta\omega_n}$
    \item $5\%$ 误差带：$\displaystyle t_s=\frac{3}{\zeta\omega_n}$
\end{itemize}

补充结论1：$\zeta\omega_n$ 是闭环极点到虚轴的距离，调节时间 $t_s$ 近似与该距离成反比；极点离虚轴越远，调节时间越短。
补充结论2：当 $\zeta=\frac{\sqrt{2}}{2}$ 时，调节时间 $t_s$ 取得最小值，该阻尼比称为**最佳阻尼比**，在快速性和平稳性之间实现最优平衡。

\subsubsection{⑤ 振荡次数 $N$}
振荡次数定义：动态过程持续时间内（$t\le t_s$），单位阶跃响应振荡的次数；将 $t\le t_s$ 内响应穿越稳态值次数的一半定义为振荡次数 $N$，表征系统平稳性。

振荡角频率 $\omega_d = \omega_n\sqrt{1-\zeta^2}$，振荡周期 $T_d=\frac{2\pi}{\omega_d}$，因此：
\[
N=\frac{t_s}{T_d}=\frac{t_s \omega_d}{2\pi}
\]
代入 $t_s$ 近似公式：
\[
N=
\begin{cases}
\displaystyle \frac{2\sqrt{1-\zeta^2}}{\pi\zeta} & \Delta=0.02 \\[4pt]
\displaystyle \frac{1.5\sqrt{1-\zeta^2}}{\pi\zeta} & \Delta=0.05
\end{cases}
\]

\subsubsection{指标汇总}
\[
\begin{cases}
t_r = \dfrac{\pi-\theta}{\omega_n\sqrt{1-\zeta^2}},\quad \theta=\arctan\dfrac{\sqrt{1-\zeta^2}}{\zeta}\\[6pt]
t_p = \dfrac{\pi}{\omega_n\sqrt{1-\zeta^2}}\\[6pt]
t_s =
\begin{cases}
\dfrac{4}{\zeta\omega_n} & \Delta=0.02\\[4pt]
\dfrac{3}{\zeta\omega_n} & \Delta=0.05
\end{cases}\\[6pt]
\sigma_p = e^{-\frac{\zeta\pi}{\sqrt{1-\zeta^2}}}\times 100\%
\end{cases}
\]

\subsection{3.3.4 二阶系统的单位脉冲、单位斜坡响应}
\subsubsection{1. 单位脉冲响应}
系统脉冲传递函数：
\[
Y(s)=\frac{\omega_n^2}{s^2+2\zeta\omega_n s+\omega_n^2}
\]
\begin{enumerate}
    \item \textbf{无阻尼（$\zeta=0$）}
    \[
    k(t)=\omega_n \sin\omega_n t,\quad t\ge0
    \]
    \item \textbf{欠阻尼（$0<\zeta<1$）}
    \[
    k(t)=\frac{\omega_n}{\sqrt{1-\zeta^2}}e^{-\zeta\omega_n t}\sin\omega_d t,\quad t\ge0
    \]
    \item \textbf{临界阻尼（$\zeta=1$）}
    \[
    k(t)=\omega_n^2 t e^{-\omega_n t},\quad t\ge0
    \]
    \item \textbf{过阻尼（$\zeta>1$）}
    \[
    k(t)=\frac{\omega_n}{2\sqrt{\zeta^2-1}}\left[e^{-(\zeta-\sqrt{\zeta^2-1})\omega_n t}-e^{-(\zeta+\sqrt{\zeta^2-1})\omega_n t}\right],\quad t\ge0
    \]
\end{enumerate}

\paragraph{脉冲响应与阶跃响应关系}
单位脉冲函数是单位阶跃函数对时间的导数，因此**单位脉冲响应 = 单位阶跃响应对时间求导**；对阶跃响应求导即可得到脉冲响应。
脉冲响应是传递函数的拉普拉斯逆变换，和传递函数等价，完整表征控制系统全部特性。

积分关系：
\[
y(t_p)=\int_0^{t_p} k(t)dt = 1+\sigma_p
\]

\subsubsection{2. 单位斜坡响应 $r(t)=t$}
\[
Y(s)=\frac{\omega_n^2}{s^2+2\zeta\omega_n s+\omega_n^2}\cdot \frac{1}{s^2}
=\frac{1}{s^2}-\frac{2\zeta/\omega_n}{s}+\frac{2\zeta(s+\zeta\omega_n)/\omega_n+(2\zeta^2-1)}{s^2+2\zeta\omega_n s+\omega_n^2}
\]

\begin{enumerate}
    \item \textbf{欠阻尼（$0<\zeta<1$）斜坡响应}
    \[
    \begin{aligned}
    y(t) &= t-\frac{2\zeta}{\omega_n}+e^{-\zeta\omega_n t}\left( \frac{2\zeta}{\omega_n}\cos\omega_d t+\frac{2\zeta^2-1}{\omega_n\sqrt{1-\zeta^2}}\sin\omega_d t \right) \\
    &= t-\frac{2\zeta}{\omega_n}+\frac{1}{\omega_n\sqrt{1-\zeta^2}}e^{-\zeta\omega_n t}\sin\left( \omega_d t+\arctan\frac{2\zeta\sqrt{1-\zeta^2}}{2\zeta^2-1} \right),\quad t\ge0
    \end{aligned}
    \]
    \item \textbf{临界阻尼（$\zeta=1$）斜坡响应}
    \[
    y(t)=t-\frac{2}{\omega_n}+\frac{2}{\omega_n}e^{-\omega_n t}\left(1+\frac{\omega_n t}{2}\right),\quad t\ge0
    \]
    \item \textbf{过阻尼（$\zeta>1$）斜坡响应}
    \[
    \begin{aligned}
    y(t) &= t-\frac{2\zeta}{\omega_n}-\frac{2\zeta^2-1-2\zeta\sqrt{\zeta^2-1}}{2\omega_n\sqrt{\zeta^2-1}}e^{-(\zeta+\sqrt{\zeta^2-1})\omega_n t} \\
    &\quad +\frac{2\zeta^2-1+2\zeta\sqrt{\zeta^2-1}}{2\omega_n\sqrt{\zeta^2-1}}e^{-(\zeta-\sqrt{\zeta^2-1})\omega_n t},\quad t\ge0
    \end{aligned}
    \]
\end{enumerate}

\subsection{欠阻尼二阶系统极点位置与动态性能变化规律}
闭环共轭极点：
\[
s_{1,2}=-\zeta\omega_n \pm j\omega_n\sqrt{1-\zeta^2}
\]
\[
c(t) = 1-\frac{e^{-\zeta\omega_n t}}{\sqrt{1-\zeta^2}}\sin(\omega_d t+\theta)
\]

\subsubsection{1. 等实部垂线移动（实部$-\zeta\omega_n$固定）}
- 实部不变 $\Rightarrow$ 衰减系数 $e^{-\zeta\omega_n t}$ 不变，调节时间基本不变；
- 虚部 $\omega_d$ 增大 $\Rightarrow$ 振荡频率升高，上升/峰值时间缩短。

\subsubsection{2. 极点沿原点射线向外/向内移动（$\zeta$、$\theta$固定不变）}
- 极点远离原点：$\omega_n$ 增大；上升时间、峰值时间、调节时间全部缩短，响应变快；超调量不变。
- 极点靠近原点：$\omega_n$ 减小；各项时间指标全部变长，响应变慢；超调量不变。

\begin{table}[htbp]
\centering
\caption{极点沿射线移动性能变化（$\zeta$恒定）}
\begin{tabular}{ccccc}
\toprule
极点移动方向 & 自然频率$\omega_n$ & 超调量$\sigma_p$ & 上升/峰值时间 & 调节时间 \\
\midrule
远离原点 & $\uparrow$增大 & 不变 & $\downarrow$显著变快 & $\downarrow$显著变快 \\
靠近原点 & $\downarrow$减小 & 不变 & $\uparrow$显著变慢 & $\uparrow$显著变慢 \\
\bottomrule
\end{tabular}
\end{table}

\subsubsection{3. 极点沿同心圆（$\omega_n$固定）圆弧移动}
- 向虚轴靠近：$\theta$ 变大，$\zeta=\cos\theta$ 减小；超调量剧增，上升时间变快，振荡次数增多，调节时间变长；
- 向负实轴靠近：$\theta$ 变小，$\zeta=\cos\theta$ 增大；超调量减小，上升时间变慢，振荡次数减少，调节时间先减小后增大（存在最小值）。

\begin{table}[htbp]
\centering
\caption{极点沿圆弧移动性能变化（$\omega_n$恒定）}
\begin{tabular}{cccccc}
\toprule
极点移动方向 & 阻尼比$\zeta$ & 超调量$\sigma_p$ & 上升时间 & 调节时间 & 振荡次数 \\
\midrule
靠近虚轴 & $\downarrow$减小 & $\uparrow$剧增 & $\downarrow$变快 & $\uparrow$振荡久、变慢 & $\uparrow$增多 \\
靠近负实轴 & $\uparrow$增大 & $\downarrow$减小 & $\uparrow$变慢 & 先$\downarrow$后$\uparrow$（极小值） & $\downarrow$减少 \\
\bottomrule
\end{tabular}
\end{table}

\bigskip

% ===== 恢复全局默认间距（避免影响后续章节）=====
\setlength{\parskip}{0pt plus 1pt}
\setlist{leftmargin=2em, itemsep=0.2em, parsep=0pt}
\renewcommand{\arraystretch}{1}
\setlength{\tabcolsep}{6pt}
% ============================================